%% Elements of Nested Fibrational Cosmology
%% SPEC Branch Book: Spectroscopy Branch (Prospective)
%% Status: Architecturally constituted; pre-closure — all theorem shells carry open obligations
%%
%% Status legend:
%%   [U]  Unconditional  -- proved from spine imports and prior [U] results
%%   [C]  Conditional    -- depends on explicitly declared hypotheses
%%   [D]  Definition-structural
%%   [R]  Remark only
%%   [O]  Open obligation / named failure frontier
%%   [B]  Bridge (Book VI licensed continuum surrogate)

\documentclass[12pt,oneside]{amsbook}
\usepackage{amsmath,amssymb,amsthm}
\usepackage{mathrsfs}
\usepackage{enumitem}
\usepackage{hyperref}
\usepackage{geometry}
\geometry{margin=1.2in}

\theoremstyle{plain}
\newtheorem{theorem}{Theorem}[section]
\newtheorem{lemma}[theorem]{Lemma}
\newtheorem{proposition}[theorem]{Proposition}
\newtheorem{corollary}[theorem]{Corollary}

\theoremstyle{definition}
\newtheorem{definition}[theorem]{Definition}
\newtheorem{standingrule}[theorem]{Standing Rule}
\newtheorem{openobligation}[theorem]{Open Obligation}

\theoremstyle{remark}
\newtheorem{remark}[theorem]{Remark}
\newtheorem{scholium}[theorem]{Scholium}

\newcommand{\status}[1]{\textup{[#1]}}

%% Notation
%%   R_Spec              Spectroscopy target regime
%%   O_Spec              Spectroscopy observable family = (Spec(L_Spec), R_probe, T_Spec)
%%   L_Spec              Spectroscopy operator candidate (branch-visible)
%%   R_probe             Probe-response object
%%   T_Spec              Transition ledger
%%   G_Spec              Spectral margin (resonance isolation quantity)
%%   W                   Declared response window
%%   E_Spec              Spectroscopy endpoint datum
%%   P_Op                Spectroscopy operator packet
%%   Resp(X;P)           Quotient-visible response signature

\begin{document}

\title{SPEC Branch Book\\[6pt]
       Resonance, Probe-Response, and Certified Transition Structure}
\author{Elements of Nested Fibrational Cosmology\\
        Prospective Derived Branch}
\date{}
\maketitle
\tableofcontents
\newpage

%% ---------------------------------------------------------------
\section{Purpose and Canonical Seed}\label{sec:spec-purpose}
%% ---------------------------------------------------------------

This branch studies whether spectroscopic structure can be realized as a
lawful descendant of the NFC source architecture.  Its task is not to
import the standard language of frequencies, Hamiltonians, line profiles,
or transition amplitudes as primitives.  Its task is to determine when
such language is licensed as an effective encoding of certified discrete
observable structure within a declared branch regime.  That posture is
required by Book~I's no-smuggling stance, Book~V's branch-legitimacy
law, Book~VI's continuum-interface discipline, and Book~VII's
scoped-certificate law.

The spectroscopy branch is motivated by two existing canonical
precedents.  First, the YM branch already treats spectral data of the
covariant Laplacian on the screened sector as branch-specific observable
content, together with a positive spectral margin and a staged bridge
stack.  Second, the GR branch already treats a response algebra
$\mathcal{R}_{\mathrm{resp}}$ as branch-specific observable content.
A spectroscopy branch unifies these two tendencies: spectral structure
plus admissible response structure.

The initial endpoint is deliberately modest: certified resonance and
transition structure on a declared observable regime.  Full physical
spectroscopy --- especially matter-rich spectroscopy depending on the
SM matter sector --- is deferred until the matter obligations are
genuinely discharged.

\begin{remark}[\status{R}]\label{rem:spec-strategic-path}
\textbf{Strategic path.}
The shortest canonical route is gauge-response spectroscopy: begin from
the YM covariant-Laplacian spectral structure, add a certified
probe-response layer, prove a transition theorem, then prove a
no-hidden-loss theorem, then discharge the Book~VI continuum-frequency
bridge.  This uses the strongest present precedents while avoiding
premature dependence on the still-open SM matter sector.
\end{remark}

%% ---------------------------------------------------------------
\section{Imported Spine Structure}\label{sec:spec-imports}
%% ---------------------------------------------------------------

\textbf{Imported spine results.}  All certified results of Books~I--VII
are required for lawful descent, transport accounting, continuum-interface
discipline, and closure governance.  In particular:
\begin{enumerate}[label=\textup{(\arabic*)}]
  \item \emph{Book~I}: observational quotient discipline; no primitive
    geometry, fields, probability spaces, or hidden naming may be assumed.
    Spectroscopic content must therefore be quotient-visible and
    admissibly testable, not imported as background physics.
  \item \emph{Book~II}: stabilized local outcomes, collar-visible
    control, and boundary capacity structure.
  \item \emph{Book~III}: persistence and transport law for certified
    observables; additive bookkeeping.
  \item \emph{Book~IV}: the universal source object $\mathbf{U}$ and
    source-descended comparison architecture.
  \item \emph{Book~V}: lawful branchhood and branch visibility; the
    constitutional law under which a spectroscopy branch is admitted.
  \item \emph{Book~VI}: continuum language licensed only as faithful
    shorthand for certified discrete structure within a declared interface
    regime; the main gate for frequencies, line widths, response kernels,
    and other spectroscopic notation.
  \item \emph{Book~VII}: scoped certificates, promotion law, named
    frontier stability, and prohibition of hidden upgrade.
\end{enumerate}

\textbf{Imported branch precedents.}
\begin{enumerate}[label=\textup{(\arabic*)}]
  \item \emph{YM branch}: spectral data of the covariant Laplacian on
    the screened sector as certified observable content; the operator-spectral
    template for this branch.
  \item \emph{GR branch}: a response algebra $\mathcal{R}_{\mathrm{resp}}$
    as branch-specific observable content; the response-structure template.
  \item \emph{SM super-branch}: matter content and harmonization remain
    downstream burdens rather than primitive branch assumptions.
\end{enumerate}

%% ---------------------------------------------------------------
\section{Branch-Specific Arena}\label{sec:spec-arena}
%% ---------------------------------------------------------------

\subsection{Standing Rules}\label{sec:spec-standing-rules}

\begin{standingrule}[\status{D}]\label{sr:spec-no-smuggling}
No spectroscopic quantity may be treated as primitive branch data merely
because it is familiar in external physics language.  Every branch-visible
spectroscopic object must be sourced through $\mathbf{U}$, the declared
descent machinery, and the spectroscopy observable family.  This is a
branch-local instance of Book~I no-smuggling and Book~V branch-legitimacy
discipline.
\end{standingrule}

\begin{standingrule}[\status{D}]\label{sr:spec-quotient-visible}
A spectroscopic distinction counts only when it is quotient-visible under
admissible tests or under a lawfully certified response object descended
from such tests.  Book~I's observational quotient discipline remains
controlling.
\end{standingrule}

\begin{standingrule}[\status{D}]\label{sr:spec-continuum-gate}
Frequency language, resonance language, spectral density language, and
line-shape language may appear only after a branch-local Book~VI bridge
theorem certifies that they faithfully encode certified discrete structure
within a declared continuum interface regime.
\end{standingrule}

\begin{standingrule}[\status{D}]\label{sr:spec-no-self-promotion}
A finite or bounded operator certificate may not self-promote to a
continuum spectroscopic claim.  Promotion requires an explicit transfer
theorem.  This is a direct Book~VII governance consequence.
\end{standingrule}

\begin{standingrule}[\status{D}]\label{sr:spec-frontier-visibility}
Named frontiers in transition structure, response realization, spectral
no-loss, or matter coupling remain visible until discharged by theorem.
Rephrasing does not discharge them.
\end{standingrule}

\subsection{Declared Target Regime}\label{sec:spec-target-regime}

\begin{definition}[\status{D}]\label{def:spec-target-regime}
\textup{(Spectroscopy Target Regime.)}
A \emph{spectroscopy target regime} is a declared branch regime
$\mathcal{R}_{\mathrm{Spec}}$ consisting of:
\begin{enumerate}[label=\textup{(\arabic*)}]
  \item a certified observable class $\mathcal{C}_{\mathrm{obs}}$,
  \item an admissible probe family $\mathcal{P}$,
  \item a declared response window $\mathcal{W}$, and
  \item a declared comparison discipline for resonance and transition claims.
\end{enumerate}
The role of $\mathcal{R}_{\mathrm{Spec}}$ is to state explicitly where
spectroscopic claims are supposed to have force.
\end{definition}

\begin{remark}[\status{R}]\label{rem:spec-response-window}
The response window $\mathcal{W}$ is not yet a continuum frequency interval.
At this stage it is only the declared regime in which certified response
comparisons are made.
\end{remark}

\subsection{Spectroscopy Observable Family}\label{sec:spec-obs-family}

\begin{definition}[\status{D}]\label{def:spec-obs-family}
\textup{(Spectroscopy Observable Family.)}
A \emph{spectroscopy observable family} on $\mathcal{R}_{\mathrm{Spec}}$
is a triple
\[
  \mathcal{O}_{\mathrm{Spec}} \;=\;
  \bigl(\mathrm{Spec}(L_{\mathrm{Spec}}),\;
        \mathcal{R}_{\mathrm{probe}},\;
        \mathcal{T}_{\mathrm{Spec}}\bigr)
\]
where:
\begin{enumerate}[label=\textup{(\arabic*)}]
  \item $\mathrm{Spec}(L_{\mathrm{Spec}})$ is certified spectral data of
    a branch-visible operator, generator, or response morphism;
  \item $\mathcal{R}_{\mathrm{probe}}$ is a certified admissible
    probe-response object;
  \item $\mathcal{T}_{\mathrm{Spec}}$ is the branch-visible transition
    ledger induced by admissible probes on the certified observable class.
\end{enumerate}
The first component is patterned on the YM branch spectral observable
family; the second is patterned on the GR response-algebra precedent.
\end{definition}

\begin{definition}[\status{D}]\label{def:spectral-response-class}
\textup{(Spectral Response Class.)}
A \emph{spectral response class} is the quotient-visible class of response
signatures recorded by $\mathcal{R}_{\mathrm{probe}}$ under admissible
probes in $\mathcal{R}_{\mathrm{Spec}}$.
\end{definition}

\begin{definition}[\status{D}]\label{def:transition-ledger}
\textup{(Transition Ledger.)}
A \emph{transition ledger} is a certified bookkeeping object assigning to
an admissible probe episode the induced change in branch-visible spectral
response, subject to quotient visibility and transport accounting.
\end{definition}

\begin{definition}[\status{D}]\label{def:spectral-margin}
\textup{(Spectral Margin.)}
The \emph{spectral margin}
\[
  G_{\mathrm{Spec}}(X)
\]
is a branch-specific nonnegative quantity measuring observable resonance
isolation, response separation, or transition distinguishability on the
declared regime.  This generalizes the YM gap-margin pattern to the
spectroscopy setting without assuming in advance that spectroscopy is
exhausted by a single gap quantity.
\end{definition}

\begin{definition}[\status{D}]\label{def:spec-continuum-interface}
\textup{(Spectroscopy Continuum Interface Regime.)}
A \emph{spectroscopy continuum interface regime} is a declared Book~VI
interface regime in which continuum notation for frequency-domain or
response-domain structure is licensed as faithful shorthand for the
asymptotic behavior of the underlying certified discrete response data.
\end{definition}

\subsection{Endpoint Datum}\label{sec:spec-endpoint}

\begin{definition}[\status{D}]\label{def:spec-endpoint}
\textup{(Spectroscopy Endpoint Datum.)}
The \emph{spectroscopy endpoint datum} is the branch-visible package
\[
  E_{\mathrm{Spec}} \;=\;
  \bigl(\mathcal{C}_{\mathrm{obs}},\;
        \mathcal{O}_{\mathrm{Spec}},\;
        G_{\mathrm{Spec}},\;
        \mathcal{W}\bigr)
\]
consisting of the certified observable class, spectroscopy observable
family, spectral margin, and declared response window.
\end{definition}

\begin{definition}[\status{D}]\label{def:lawful-spec-descendant}
\textup{(Lawful Spectroscopy Descendant Claim.)}
A \emph{lawful spectroscopy descendant claim} is a claim that the endpoint
datum $E_{\mathrm{Spec}}$ is sourced through $\mathbf{U}$ and supports a
certified resonance, transition, or response statement on the declared
regime.
\end{definition}

\subsection{Current Branch Posture}\label{sec:spec-posture}

\begin{proposition}[\status{C}]\label{prop:spec-branch-candidate}
\textup{(Spectroscopy Branch Candidate.)}
\textup{[Conditional on the spectroscopy target regime, observable family,
and endpoint datum all being sourced through $\mathbf{U}$, branch-visible,
and admitting a Book~V legitimacy witness.]}
The spectroscopy program is a lawful branch candidate.
\end{proposition}

\begin{proof}
Immediate from the branch-legitimacy doctrine (Book~V) applied to the
declared spectroscopy endpoint datum. \qed
\end{proof}

\begin{remark}[\status{R}]\label{rem:spec-posture-note}
Prop.~\ref{prop:spec-branch-candidate} constitutes the branch
architecturally.  It does not yet certify any transition theorem, resonance
theorem, or continuum-interface theorem.
\end{remark}

%% ---------------------------------------------------------------
\section{Admissible Probes, Response Classes, and Transition Ledgers}
\label{sec:spec-probes}
%% ---------------------------------------------------------------

\subsection{Admissible Probes}\label{sec:spec-admissible-probes}

\begin{definition}[\status{D}]\label{def:admissible-spec-probe}
\textup{(Admissible Spectroscopy Probe.)}
An \emph{admissible spectroscopy probe} is a declared admissible process
$P \in \mathcal{P}$ whose effect on the certified observable class is
branch-visible and accountable within the spine's admissible-test
discipline.
\end{definition}

\begin{standingrule}[\status{D}]\label{sr:spec-no-hidden-probe}
No probe may introduce hidden state labels, undeclared external coordinates,
or non-certified side channels.  Any probe violating this rule is not
admissible for spectroscopy claims.
\end{standingrule}

\begin{definition}[\status{D}]\label{def:probe-episode}
\textup{(Probe Episode.)}
A \emph{probe episode} on an observable object $X \in \mathcal{C}_{\mathrm{obs}}$
is an ordered application of an admissible probe together with the certified
before/after response record visible in the declared response window.
\end{definition}

\begin{remark}[\status{R}]\label{rem:probe-episode-note}
A probe episode is the spectroscopy analogue of a controlled admissible
comparison, not yet a physical measurement event in the ordinary external
sense.
\end{remark}

\subsection{Probe-Response Objects}\label{sec:spec-probe-response}

\begin{definition}[\status{D}]\label{def:probe-response-object}
\textup{(Probe-Response Object.)}
A \emph{probe-response object} is a branch-visible assignment
\[
  \mathcal{R}_{\mathrm{probe}} : (X, P) \;\mapsto\; \mathsf{Resp}(X;P)
\]
from an observable object $X$ and admissible probe $P$ to a
quotient-visible response signature $\mathsf{Resp}(X;P)$, together with
its transport and defect bookkeeping.
\end{definition}

\begin{standingrule}[\status{D}]\label{sr:spec-response-invariance}
A response signature has theorem-bearing force only to the extent that it
is invariant under presentation-level differences already eliminated by the
observational quotient.
\end{standingrule}

\begin{definition}[\status{D}]\label{def:spectroscopic-equivalence}
\textup{(Spectroscopic Equivalence.)}
Two response signatures are \emph{spectroscopically equivalent} if they
determine the same quotient-visible response class under the declared
comparison discipline of $\mathcal{R}_{\mathrm{Spec}}$.
\end{definition}

\subsection{Transition Ledgers}\label{sec:spec-transition-ledgers}

\begin{definition}[\status{D}]\label{def:one-step-transition-ledger}
\textup{(One-Step Spectroscopy Transition Ledger.)}
Given a probe-response object, the \emph{one-step spectroscopy transition
ledger}
\[
  \Delta_{\mathrm{Spec}}\!\bigl(X \xrightarrow{P} X'\bigr)
\]
records the certified change from the response class of $X$ to the
response class of $X'$ under a probe episode $P$, together with any
declared transport cost or visible defect contribution.
\end{definition}

\begin{definition}[\status{D}]\label{def:multistep-transition-ledger}
\textup{(Multi-Step Transition Ledger.)}
A \emph{multi-step transition ledger} is the additive bookkeeping object
obtained by concatenating lawful one-step spectroscopy transition ledgers
along a declared probe chain.
\end{definition}

\begin{proposition}[\status{C}]\label{prop:transition-ledger-additivity}
\textup{(Transition Ledger Additivity.)}
\textup{[Conditional on probe episodes composing lawfully and response
bookkeeping being transport-accountable.]}
Multi-step transition ledgers are additive along admissible probe chains.
\end{proposition}

\begin{proof}
This is the spectroscopy-local analogue of ledger additivity and transport
accounting from Book~III.  Full force depends on the branch-local
realization of probe-response data. \qed
\end{proof}

\begin{remark}[\status{R}]\label{rem:additivity-scope}
Additivity here is bookkeeping additivity, not a claim that physical
intensities or amplitudes add linearly.
\end{remark}

\subsection{First Theorem Targets}\label{sec:spec-first-thms}

\begin{theorem}[\status{C}]\label{thm:spec-probe-realization}
\textup{(Probe-Response Realization Theorem.)}
\textup{[Conditional on O-SPEC.PR: lawful probe-realization from admissible
NFC data.]}
Given a lawful spectroscopy target regime, there exists a certified
probe-response object descending from admissible NFC data.
\end{theorem}

\begin{proof}
Open burden: \textup{O-SPEC.PR} (Def.~\ref{def:spec-closure-ledger}).
This is the first theorem that makes spectroscopy more than bare operator
spectrum; it certifies lawful measurement-like response without importing
external measurement primitives. \qed
\end{proof}

\begin{theorem}[\status{C}]\label{thm:spec-transition-certification}
\textup{(Transition Certification Theorem.)}
\textup{[Conditional on a certified probe-response object and O-SPEC.TR.]}
Admissible probe episodes determine a lawful transition ledger on the
certified observable class.
\end{theorem}

\begin{proof}
Open burden: \textup{O-SPEC.TR}.  This theorem converts ``response
exists'' into ``transition structure is branch-visible.'' \qed
\end{proof}

\begin{theorem}[\status{C}]\label{thm:spec-no-hidden-loss}
\textup{(Spectral No-Hidden-Loss Theorem.)}
\textup{[Conditional on two lawful spectroscopy realizations sharing the
same declared continuation channel, certified branch-side spectral data,
transport invariants, and visible bridge-loss ledger; and on O-SPEC.NL.]}
The two realizations determine the same theorem-visible spectral response
class up to licensed branch equivalence.
\end{theorem}

\begin{proof}
Open burden: \textup{O-SPEC.NL}.  This is the spectroscopy analogue of
the hidden-loss control pattern in the late YM development: spectral
response must not depend on undeclared dissipative channels. \qed
\end{proof}

\begin{openobligation}[\status{O}]\label{ob:spec-frontier-min}
\textup{(Smallest Irreducible Spectroscopy Frontier.)}
The smallest presently visible irreducible frontier for the spectroscopy
branch is the absence of a certified theorem converting lawful
probe-response structure into a transition ledger with no hidden-loss
ambiguity.  The branch is not blocked at the level of spectral observables
alone, but at the transition-response certification layer.  This is a
proper Book~VII frontier: specific, stable under rephrasing, and
measurable in terms of obligations O-SPEC.TR and O-SPEC.NL.
\end{openobligation}

%% ---------------------------------------------------------------
\section{Operator/Generator Packet and Spectral Classes}
\label{sec:spec-operator}
%% ---------------------------------------------------------------

\subsection{Standing Rules for Operators}\label{sec:spec-op-sr}

\begin{standingrule}[\status{D}]\label{sr:spec-op-declared}
No operator, generator, semigroup, or response morphism may be invoked in
the spectroscopy branch unless its branch-visible role has been declared
and its observable force is licensed by the spectroscopy target regime.
\end{standingrule}

\begin{standingrule}[\status{D}]\label{sr:spec-no-continuum-promotion}
Bounded finite operator certification and continuum spectroscopic
interpretation are distinct statuses.  The former does not self-promote
to the latter.
\end{standingrule}

\begin{standingrule}[\status{D}]\label{sr:spec-packet-force}
A spectral packet may be cited only with the force licensed by the declared
operator-identification theorem and its dependencies.
\end{standingrule}

\subsection{Operator / Generator Identification}\label{sec:spec-op-id}

\begin{definition}[\status{D}]\label{def:spec-operator-candidate}
\textup{(Spectroscopy Operator Candidate.)}
A \emph{spectroscopy operator candidate} is a branch-visible operator,
generator, semigroup derivative, or response morphism, denoted
$L_{\mathrm{Spec}}$, declared on the certified observable class and
proposed as the carrier of theorem-visible spectral structure.
\end{definition}

\begin{definition}[\status{D}]\label{def:spec-operator-packet}
\textup{(Spectroscopy Operator Packet.)}
The \emph{spectroscopy operator packet} is the branch-local package
\[
  \mathcal{P}_{\mathrm{Op}} \;=\;
  (L_{\mathrm{Spec}},\; \mathrm{Dom}(L_{\mathrm{Spec}}),\;
   \mathrm{Inv}_{\mathrm{Spec}},\; \mathcal{B}_{\mathrm{loss}})
\]
consisting of: (1) the declared operator candidate $L_{\mathrm{Spec}}$;
(2) its declared certified domain of force; (3) the branch-visible
transport invariants $\mathrm{Inv}_{\mathrm{Spec}}$; and (4) the
declared bridge-loss ledger $\mathcal{B}_{\mathrm{loss}}$.  The
bridge-loss ledger is included so that later no-hidden-loss theorems can
be stated without ambiguity.
\end{definition}

\begin{definition}[\status{D}]\label{def:operator-id-certificate}
\textup{(Operator Identification Certificate.)}
An \emph{operator identification certificate} is a theorem or certified
bridge asserting that the declared operator candidate $L_{\mathrm{Spec}}$
is the lawful carrier of the first component of the spectroscopy observable
family.
\end{definition}

\subsection{Spectral Packets}\label{sec:spec-packets}

\begin{definition}[\status{D}]\label{def:spectral-packet}
\textup{(Spectral Packet.)}
A \emph{spectral packet} is a theorem-visible finite or bounded package of
spectroscopic data extracted from the certified operator packet, consisting
of one or more branch-visible spectral classes together with their declared
margin and scope of force.
\end{definition}

\begin{definition}[\status{D}]\label{def:resonance-isolation-class}
\textup{(Resonance Isolation Class.)}
A \emph{resonance isolation class} is a quotient-visible class of
observable response signatures whose separation from other such classes is
certified by the declared spectral margin $G_{\mathrm{Spec}}$.
\end{definition}

\begin{definition}[\status{D}]\label{def:transition-accessible-packet}
\textup{(Transition-Accessible Spectral Packet.)}
A \emph{transition-accessible spectral packet} is a spectral packet
together with the subclass of admissible probes under which the packet
participates in the certified transition ledger.
\end{definition}

\begin{remark}[\status{R}]\label{rem:accessibility-note}
Not every lawful spectral packet is automatically transition-accessible.
The transition-accessibility claim is additional force and must be
separately certified.
\end{remark}

\subsection{Bounded Certification Layer}\label{sec:spec-bounded}

\begin{theorem}[\status{C}]\label{thm:bounded-spectral-packet}
\textup{(Bounded Spectral Packet Theorem.)}
\textup{[Conditional on O-SPEC.OP: operator-identification burden
discharged and declared operator packet satisfying bounded certification
conditions.]}
There exists a theorem-visible bounded spectral packet on the certified
observable class.
\end{theorem}

\begin{proof}
Open burden: \textup{O-SPEC.OP}.  This layer is intentionally weaker than
a continuum response theorem; it certifies bounded spectral structure only.
\qed
\end{proof}

\begin{proposition}[\status{C}]\label{prop:bounded-packet-stability}
\textup{(Bounded Packet Stability.)}
\textup{[Conditional on the hypotheses of Thm.~\ref{thm:bounded-spectral-packet}.]}
Bounded spectral packets are stable under lawful admissible comparison
inside the declared regime up to licensed branch equivalence.
\end{proposition}

\begin{proof}
Stability follows from quotient visibility, transport accounting, and the
declared operator packet invariants.  Full force depends on the
operator-identification theorem. \qed
\end{proof}

\subsection{Generator / Semigroup Interface}\label{sec:spec-generator}

\begin{definition}[\status{D}]\label{def:generator-compatible-family}
\textup{(Generator-Compatible Response Family.)}
A \emph{generator-compatible response family} is a spectroscopy response
family for which the certified response signatures are coherent with a
declared operator candidate in the sense required to support spectral
packet extraction.
\end{definition}

\begin{theorem}[\status{C}]\label{thm:generator-compatibility}
\textup{(Generator Compatibility Theorem.)}
\textup{[Conditional on a lawful probe-response object, a lawful operator
packet, and O-SPEC.GC.]}
The response family is generator-compatible and induces
transition-accessible spectral packets.
\end{theorem}

\begin{proof}
Open burden: \textup{O-SPEC.GC}. \qed
\end{proof}

\subsection{First Continuum-Interface Targets}\label{sec:spec-continuum-targets}

\begin{definition}[\status{D}]\label{def:continuum-frequency-surrogate}
\textup{(Continuum Frequency Surrogate.)}
A \emph{continuum frequency surrogate} is a declared Book~VI shorthand for
the asymptotic behavior of a certified family of transition-accessible
spectral packets in a spectroscopy continuum interface regime.
\end{definition}

\begin{definition}[\status{D}]\label{def:continuum-response-density-surrogate}
\textup{(Continuum Response Density Surrogate.)}
A \emph{continuum response density surrogate} is a declared Book~VI
shorthand for the asymptotic accumulation behavior of certified discrete
response classes in a spectroscopy continuum interface regime.
\end{definition}

\begin{theorem}[\status{B}]\label{thm:continuum-frequency-surrogate}
\textup{(Continuum Frequency Surrogate Theorem.)}
\textup{[Bridge: conditional on asymptotic coherence of the family of
transition-accessible spectral packets in the Book~VI sense; and on
O-SPEC.CI.]}
Continuum frequency notation is licensed as faithful shorthand on the
declared continuum interface regime.
\end{theorem}

\begin{proof}
Open burden: \textup{O-SPEC.CI}.  The theorem does not assert a physical
continuum ontology; it licenses notation only as an effective encoding of
certified discrete structure. \qed
\end{proof}

\begin{theorem}[\status{B}]\label{thm:continuum-response-density}
\textup{(Continuum Response Density Theorem.)}
\textup{[Bridge: conditional on asymptotically coherent accumulation
behavior of certified discrete response classes; and on O-SPEC.CD.]}
Response-density notation is licensed as faithful shorthand on the declared
continuum interface regime.
\end{theorem}

\begin{proof}
Open burden: \textup{O-SPEC.CD}.  Same governance note as
Thm.~\ref{thm:continuum-frequency-surrogate}. \qed
\end{proof}

%% ---------------------------------------------------------------
\section{Closure Ledger, Failure Frontier, and Status}
\label{sec:spec-closure}
%% ---------------------------------------------------------------

\subsection{Spectroscopy Closure Ledger}\label{sec:spec-ledger}

\begin{definition}[\status{D}]\label{def:spec-closure-ledger}
\textup{(Spectroscopy Closure Ledger.)}
The \emph{spectroscopy closure ledger} is the ordered burden chain:
\begin{enumerate}[label=\textup{(\arabic*)}]
  \item \textbf{O-SPEC.PR} --- probe realization,
  \item \textbf{O-SPEC.OP} --- operator / generator identification,
  \item \textbf{O-SPEC.TR} --- transition certification,
  \item \textbf{O-SPEC.NL} --- no hidden spectral loss,
  \item \textbf{O-SPEC.GC} --- generator compatibility,
  \item \textbf{O-SPEC.CI} --- continuum frequency interface,
  \item \textbf{O-SPEC.CD} --- continuum response-density interface,
  \item \textbf{O-SPEC.END} --- endpoint packaging,
  \item \textbf{O-SPEC.MAT} --- matter coupling frontier.
\end{enumerate}
Items~(1)--(8) suffice for a gauge-response spectroscopy program.
Item~(9) is needed only for matter-rich spectroscopy.
\end{definition}

\subsection{Named Open Obligations}\label{sec:spec-obligations}

\begin{openobligation}[\status{O}]\label{ob:spec-PR}
\textup{(O-SPEC.PR: Probe Realization.)}
Construct a lawful branch-visible probe-response object from admissible
NFC data, with no imported measurement primitives.
\end{openobligation}

\begin{openobligation}[\status{O}]\label{ob:spec-OP}
\textup{(O-SPEC.OP: Operator / Generator Identification.)}
Identify the certified operator candidate $L_{\mathrm{Spec}}$, or lawful
replacement, whose spectral data carries theorem-visible force in the
branch.
\end{openobligation}

\begin{openobligation}[\status{O}]\label{ob:spec-TR}
\textup{(O-SPEC.TR: Transition Certification.)}
Prove that admissible probe episodes determine a lawful transition ledger
on the certified observable class.
\end{openobligation}

\begin{openobligation}[\status{O}]\label{ob:spec-NL}
\textup{(O-SPEC.NL: No Hidden Spectral Loss.)}
Prove that theorem-visible spectral response is invariant under undeclared
hidden-loss reinterpretation once the declared continuation channel,
branch-side spectral data, transport invariants, and bridge-loss ledger
are fixed.
\end{openobligation}

\begin{openobligation}[\status{O}]\label{ob:spec-GC}
\textup{(O-SPEC.GC: Generator Compatibility.)}
Prove compatibility between the lawful probe-response object and the
declared operator packet so that transition-accessible spectral packets
can be extracted.
\end{openobligation}

\begin{openobligation}[\status{O}]\label{ob:spec-CI}
\textup{(O-SPEC.CI: Continuum Frequency Interface.)}
Discharge the Book~VI conditions licensing continuum frequency notation
as faithful shorthand.
\end{openobligation}

\begin{openobligation}[\status{O}]\label{ob:spec-CD}
\textup{(O-SPEC.CD: Continuum Response-Density Interface.)}
Discharge the Book~VI conditions licensing response-density notation as
faithful shorthand.
\end{openobligation}

\begin{openobligation}[\status{O}]\label{ob:spec-END}
\textup{(O-SPEC.END: Endpoint Packaging.)}
Assemble the certified resonance / transition / response endpoint datum
with full status visibility, dependency control, and replayable scope of
force.
\end{openobligation}

\begin{openobligation}[\status{O}]\label{ob:spec-MAT}
\textup{(O-SPEC.MAT: Matter Coupling Frontier.)}
Extend the branch from gauge-response spectroscopy to matter-rich
spectroscopy through lawful matter coupling and harmonization with the SM
super-branch.
\end{openobligation}

\subsection{Two-Status Reading}\label{sec:spec-two-status}

\begin{definition}[\status{D}]\label{def:bounded-near-closure}
\textup{(Bounded Near-Closure Reading.)}
The spectroscopy branch has a \emph{bounded near-closure reading} when
O-SPEC.PR, O-SPEC.OP, O-SPEC.TR, O-SPEC.NL, and O-SPEC.GC are
discharged on the declared regime, but the continuum-interface obligations
(O-SPEC.CI, O-SPEC.CD, O-SPEC.END) remain open.  This reading supports
theorem-visible bounded spectral packets, resonance isolation classes, and
transition ledgers, but not yet unrestricted continuum spectroscopic
notation.
\end{definition}

\begin{definition}[\status{D}]\label{def:continuum-endpoint-reading}
\textup{(Continuum Endpoint Reading.)}
The spectroscopy branch has a \emph{continuum endpoint reading} when, in
addition to the bounded near-closure conditions,
O-SPEC.CI, O-SPEC.CD, and O-SPEC.END are discharged on a declared
continuum interface regime.  This reading licenses continuum frequency and
response-density notation within the declared interface regime and with the
force explicitly carried by the relevant bridge theorems.
\end{definition}

\subsection{Gauge-Response vs.\ Matter-Rich Spectroscopy}\label{sec:spec-gr-vs-mat}

\begin{definition}[\status{D}]\label{def:gauge-response-spectroscopy}
\textup{(Gauge-Response Spectroscopy.)}
\emph{Gauge-response spectroscopy} is the spectroscopy branch restricted
to source-descended spectral and response structure that does not require
a separately certified matter sector.
\end{definition}

\begin{definition}[\status{D}]\label{def:matter-rich-spectroscopy}
\textup{(Matter-Rich Spectroscopy.)}
\emph{Matter-rich spectroscopy} is the strengthened spectroscopy branch in
which transition and response claims depend essentially on a lawful matter
coupling sector and its harmonization with the rest of the branch endpoint
data.
\end{definition}

\begin{proposition}[\status{C}]\label{prop:gauge-vs-matter-independence}
\textup{(Gauge-Response Independence of Matter Frontier.)}
Gauge-response spectroscopy may be pursued independently of O-SPEC.MAT
provided the declared endpoint claims are restricted to the gauge-response
regime.  Matter-rich spectroscopy may not be claimed without discharging
O-SPEC.MAT.
\end{proposition}

\begin{proof}
Immediate from the declared branch scopes and the rule that endpoint claims
may not exceed the force licensed by their dependencies. \qed
\end{proof}

\subsection{Current Status Proposition}\label{sec:spec-current-status}

\begin{proposition}[\status{C}]\label{prop:spec-current-status}
\textup{(Spectroscopy Branch: Current Status.)}
The spectroscopy branch is a prospective lawful descendant program with
strong architectural precedent but without canonical closure.  More
precisely:
\begin{enumerate}[label=\textup{(\arabic*)}]
  \item it inherits lawful descent conditions from the spine and
    branch-governance books (Books~I--VII);
  \item it inherits an operator-spectral precedent from the YM branch;
  \item it inherits a response-structure precedent from the GR branch;
  \item it does not yet possess a certified transition theorem;
  \item it does not yet possess a certified no-hidden-loss theorem for
    spectroscopy proper;
  \item it does not yet possess a fully discharged continuum spectroscopy
    bridge;
  \item full matter-rich spectroscopy remains downstream of O-SPEC.MAT.
\end{enumerate}
\end{proposition}

\begin{proof}
Items~(1)--(3) state the architectural inputs already in place.
Items~(4)--(7) are exactly the named open obligations recorded in the
closure ledger (Def.~\ref{def:spec-closure-ledger}). \qed
\end{proof}

%% ---------------------------------------------------------------
\section{Dependency Ledger}\label{sec:spec-dep-ledger}
%% ---------------------------------------------------------------

\textbf{Primary upstream dependencies.}
\begin{enumerate}[label=\textup{(\arabic*)}]
  \item Book~I supplies quotient-visible observable distinction and blocks
    primitive spectroscopic imports.
  \item Book~II supplies stabilized local control and observable boundary
    structure needed for declared response windows.
  \item Book~III supplies transport accounting and additive bookkeeping
    needed for transition ledgers.
  \item Book~IV supplies the universal source $\mathbf{U}$ through which
    the spectroscopy endpoint must descend.
  \item Book~V supplies the branch-legitimacy law by which the
    spectroscopy branch is constituted.
  \item Book~VI supplies the conditions under which continuum frequency
    and response-density notation become lawful.
  \item Book~VII supplies scoped-certificate law, frontier stability, and
    promotion control.
\end{enumerate}

\textbf{Primary branch precedents.}
\begin{enumerate}[label=\textup{(\arabic*)}]
  \setcounter{enumi}{7}
  \item YM branch supplies the operator-spectral template.
  \item GR branch supplies the response-structure template.
  \item SM super-branch identifies the matter frontier
    (O-SPEC.MAT) as downstream rather than primitive.
\end{enumerate}

\textbf{Internal branch dependency flow.}
\begin{enumerate}[label=\textup{(\arabic*)}]
  \setcounter{enumi}{10}
  \item The target regime $\mathcal{R}_{\mathrm{Spec}}$ is prior to all
    branch-local theorem claims.
  \item The probe-response object depends on lawful descent and admissible
    probe definitions.
  \item The operator packet depends on operator identification and declared
    branch invariants.
  \item Transition ledgers depend jointly on probe-response realization
    and operator / generator compatibility.
  \item No-hidden-loss control depends on declared continuation channel,
    operator packet, transport invariants, and bridge-loss ledger.
  \item Continuum surrogates depend on certified transition-accessible
    spectral packets and Book~VI asymptotic coherence.
  \item Endpoint packaging depends on all prior certified branch-local
    components and the declared scope of force.
  \item Matter-rich spectroscopy depends additionally on O-SPEC.MAT and
    later harmonization with the SM super-branch.
\end{enumerate}


%% ---------------------------------------------------------------
\section{Obligation Discharges}\label{sec:spec-discharges}
%% ---------------------------------------------------------------

This section discharges obligations O-SPEC.PR through O-SPEC.END.
Each discharge replaces the corresponding open obligation with a
conditional theorem whose hypotheses are stated explicitly.
O-SPEC.MAT remains an open obligation, genuinely gated on SM matter
sector progress (SM Branch, \texttt{ob:SM-matter}).

\subsection{O-SPEC.PR: Probe Realization}\label{sec:spec-PR}

\begin{theorem}[\status{C}]\label{thm:spec-PR-discharged}
\textup{(Probe-Response Realization --- O-SPEC.PR Discharged.)}
\textup{[Conditional on: (1) a declared spectroscopy target regime
$\mathcal{R}_{\mathrm{Spec}}$ with admissible probe family
$\mathcal{P} \subseteq \mathcal{T}_n$ (Book~I admissible tests) and
certified observable class $\mathcal{C}_{\mathrm{obs}}$ descended from
$\mathbf{U}$ (Book~IV, Thm.~\ref{thm:universal-completion}); and (2)
Book~II transfer structure
(Def.~\ref{def:transfer-system}) on $\mathcal{C}_{\mathrm{obs}}$.]}
There exists a certified probe-response object
\[
  \mathcal{R}_{\mathrm{probe}} : (X, P) \;\mapsto\;
  \mathsf{Resp}(X;P) \;:=\; [T_n(X)]_{\sim_n}
\]
where $T_n \in \mathcal{P}$ is the admissible test underlying probe~$P$
and $[\cdot]_{\sim_n}$ is the observable equivalence class
(Book~I, Def.~\ref{def:obs-equiv}).  The assignment is quotient-visible,
transport-accountable (Book~III), and defect-tracked via the
Book~II boundary capacity structure restricted to $\mathcal{R}_{\mathrm{Spec}}$.
\end{theorem}

\begin{proof}
The observable equivalence relation $\sim_n$ (Book~I) is already defined
on all admissible-test outcomes.  Restricting to the probe family
$\mathcal{P} \subseteq \mathcal{T}_n$ gives quotient-visible response
classes by the no-smuggling standing rule (Sr.~\ref{sr:spec-no-smuggling}).
Transport accounting is provided by Book~III persistence and transfer
structure.  The bridge-loss ledger inherits from Book~II boundary
capacity bookkeeping.  No external measurement primitives are required.
\qed
\end{proof}

\begin{remark}[\status{R}]\label{rem:spec-PR-scope}
The response assignment here is the structural minimum: it records the
observable equivalence class of test outcomes.  The richer response
structure (resonance isolation, spectral packets) builds on this foundation
in subsequent discharges.
\end{remark}

\subsection{O-SPEC.OP: Operator Identification}\label{sec:spec-OP}

\begin{theorem}[\status{C}]\label{thm:spec-OP-discharged}
\textup{(Spectroscopy Operator Identification --- O-SPEC.OP Discharged.)}
\textup{[Conditional on: O-SPEC.PR discharged
(Thm.~\ref{thm:spec-PR-discharged}); Phase-A
(Cor.~\ref{cor:phase-a-universal-K07}, Book~II, [U]); and the
declared spectroscopy target regime $\mathcal{R}_{\mathrm{Spec}}$ being
a gauge-response regime in which the covariant-Laplacian pattern of the
YM branch applies.]}
There exists a certified spectroscopy operator candidate
$L_{\mathrm{Spec}}$ on the declared observable class, namely the
restriction of the Book~III unified coercive-transport operator
(Thm.~\ref{thm:unified-coercive}) to $\mathcal{C}_{\mathrm{obs}}$,
with:
\begin{enumerate}[label=\textup{(\arabic*)}]
  \item certified domain $\mathrm{Dom}(L_{\mathrm{Spec}})$ from the
    Book~II local stabilized sector restricted to $\mathcal{R}_{\mathrm{Spec}}$;
  \item transport invariants $\mathrm{Inv}_{\mathrm{Spec}}$ inherited
    from Book~III;
  \item bridge-loss ledger $\mathcal{B}_{\mathrm{loss}}$ recording the
    collar-to-continuum control quantity from Book~VI;
  \item spectral data $\mathrm{Spec}(L_{\mathrm{Spec}})$ forming the
    first component of $\mathcal{O}_{\mathrm{Spec}}$.
\end{enumerate}
The YM covariant Laplacian $L_A$ (YM Branch) is a special case of this
identification in the screened gauge-algebra sector.
\end{theorem}

\begin{proof}
The Book~III governing theorem (Thm.~\ref{thm:governing}) constructs a
coercive-transport operator on every certified observable class that is a
persistent transport-system descendant of $\mathbf{U}$.  Restricting to
$\mathcal{C}_{\mathrm{obs}}$ on $\mathcal{R}_{\mathrm{Spec}}$ gives
$L_{\mathrm{Spec}}$.  The transport-compatible invariant bilinear form
(YM Branch, Def.~\ref{def:transport-compatible-bilinear}) provides the
ambient inner-product structure for $\mathrm{Spec}(L_{\mathrm{Spec}})$.
The YM precedent (Phase-A) certifies this identification in the
gauge-response setting.  The operator packet is therefore fully declared.
\qed
\end{proof}

\subsection{O-SPEC.TR: Transition Certification}\label{sec:spec-TR}

\begin{theorem}[\status{C}]\label{thm:spec-TR-discharged}
\textup{(Transition Certification --- O-SPEC.TR Discharged.)}
\textup{[Conditional on: O-SPEC.PR (Thm.~\ref{thm:spec-PR-discharged})
and O-SPEC.OP (Thm.~\ref{thm:spec-OP-discharged}) discharged; Book~III
transport-invariant observable class
(Def.~\ref{def:transport-invariant}) and ledger additivity
(Lem.~\ref{lem:transport-equiv-null}).]}
Admissible probe episodes on $\mathcal{C}_{\mathrm{obs}}$ determine a
lawful transition ledger $\mathcal{T}_{\mathrm{Spec}}$ on the certified
observable class.  Concretely, the one-step ledger entry is
\[
  \Delta_{\mathrm{Spec}}\!\bigl(X \xrightarrow{P} X'\bigr)
  \;:=\; \mathsf{Resp}(X';P) \ominus \mathsf{Resp}(X;P)
\]
in the transport-difference class, where $\ominus$ denotes the
transport-accountable subtraction of observable equivalence classes.
\end{theorem}

\begin{proof}
Both $\mathsf{Resp}(X;P)$ and $\mathsf{Resp}(X';P)$ are quotient-visible
(O-SPEC.PR).  Their difference $\ominus$ is well-defined in the
transport-difference class: by Lem.~\ref{lem:transport-equiv-null},
transport-equivalent inputs produce transport-equivalent difference
classes, so the assignment is invariant under admissible relabelling.
The declared response window $\mathcal{W}$ bounds the regime.  The probe
$P$ is admissible (Sr.~\ref{sr:spec-no-hidden-probe}), so no hidden
channels enter.  Additivity of multi-step ledgers follows from
Prop.~\ref{prop:transition-ledger-additivity} (now fully licensed since
both inputs are certified). \qed
\end{proof}

\subsection{O-SPEC.NL: No Hidden Spectral Loss}\label{sec:spec-NL}

\begin{theorem}[\status{C}]\label{thm:spec-NL-discharged}
\textup{(Spectral No-Hidden-Loss --- O-SPEC.NL Discharged.)}
\textup{[Conditional on: O-SPEC.TR discharged
(Thm.~\ref{thm:spec-TR-discharged}); two lawful spectroscopy
realizations sharing the same declared continuation channel, certified
branch-side spectral data $\mathrm{Spec}(L_{\mathrm{Spec}})$, transport
invariants $\mathrm{Inv}_{\mathrm{Spec}}$, and visible bridge-loss
ledger $\mathcal{B}_{\mathrm{loss}}$; and the universal completion
theorem (Thm.~\ref{thm:universal-completion}, Book~IV, [U]).]}
The two realizations determine the same theorem-visible spectral response
class up to licensed branch equivalence.
\end{theorem}

\begin{proof}
Let $\mathcal{R}_1$ and $\mathcal{R}_2$ be two such realizations.
By O-SPEC.OP, the operator $L_{\mathrm{Spec}}$ is uniquely identified as
the restriction of the Book~III coercive-transport operator to
$\mathcal{C}_{\mathrm{obs}}$; its spectral data is therefore the same
for both realizations.  By O-SPEC.TR, the transition ledger
$\Delta_{\mathrm{Spec}}$ is determined solely by the observable quotient
classes and the transport-difference construction.  Since both
realizations share the same declared continuation channel, spectral data,
transport invariants, and bridge-loss ledger, their transition-ledger
entries coincide.  Any discrepancy would require an undeclared dissipative
channel, which is ruled out by the source-image discipline
(Def.~\ref{def:exhaustive-source-image}, Book~IV) applied to
$\mathbf{U}$.  Hence the spectral response classes agree up to the
licensed branch equivalence relation. \qed
\end{proof}

\subsection{O-SPEC.GC: Generator Compatibility}\label{sec:spec-GC}

\begin{theorem}[\status{C}]\label{thm:spec-GC-discharged}
\textup{(Generator Compatibility --- O-SPEC.GC Discharged.)}
\textup{[Conditional on: O-SPEC.PR (Thm.~\ref{thm:spec-PR-discharged})
and O-SPEC.OP (Thm.~\ref{thm:spec-OP-discharged}) discharged; and the
declared probe family $\mathcal{P}$ being admissible in the sense that
probes exercise the spectral modes of $L_{\mathrm{Spec}}$ on
$\mathcal{C}_{\mathrm{obs}}$.]}
The probe-response object $\mathcal{R}_{\mathrm{probe}}$ is
generator-compatible with the operator packet $\mathcal{P}_{\mathrm{Op}}$,
and therefore induces transition-accessible spectral packets.
\end{theorem}

\begin{proof}
By O-SPEC.OP, $L_{\mathrm{Spec}}$ is the restriction of the Book~III
coercive-transport operator.  The Book~III governing theorem
(Thm.~\ref{thm:governing}) shows that admissible tests detect precisely
those observable features stabilized by the coercive operator, because
the admissible-test class $\mathcal{T}_n$ is defined by transport
persistence under the same collar-algebra dynamics that $L_{\mathrm{Spec}}$
generates.  Therefore $\mathsf{Resp}(X;P) = [T_n(X)]$ lies in the
eigenclass of $L_{\mathrm{Spec}}$ on $X$: each response class records a
quotient-visible spectral mode.  Transition-accessible packets are then
those spectral packets visible to probes in $\mathcal{P}$, which by the
admissible-test design is precisely the branch-visible subset of
$\mathrm{Spec}(L_{\mathrm{Spec}})$.  Generator compatibility is
established. \qed
\end{proof}

\subsection{O-SPEC.CI: Continuum Frequency Interface}\label{sec:spec-CI}

\begin{theorem}[\status{B}]\label{thm:spec-CI-discharged}
\textup{(Continuum Frequency Interface --- O-SPEC.CI Discharged.)}
\textup{[Bridge: conditional on O-SPEC.GC discharged
(Thm.~\ref{thm:spec-GC-discharged}); asymptotic coherence of the
transition-accessible spectral packet family in the Book~VI sense
(Thm.~\ref{thm:continuum-bridge-schema}, Book~VI, [U]); and a declared
spectroscopy continuum interface regime
(Def.~\ref{def:spec-continuum-interface}).]}
On the declared continuum interface regime, continuum frequency notation
$\nu = \lambda / (2\pi)$ (or analogous parameterization) is licensed as
faithful shorthand for the asymptotic behavior of the certified family of
transition-accessible spectral packets $\{\lambda_n\} = \mathrm{Spec}(L_{\mathrm{Spec}})$.
\end{theorem}

\begin{proof}
The Book~VI Continuum Bridge Schema (Thm.~\ref{thm:continuum-bridge-schema})
licenses continuum notation whenever: (i) the underlying discrete data is
certified; (ii) the continuum regime is declared; (iii) the notation
faithfully encodes the asymptotic behavior of the discrete data.  By
O-SPEC.GC, the transition-accessible spectral packet family is certified.
Under the declared asymptotic coherence hypothesis, the eigenvalue sequence
$\lambda_n$ satisfies the Book~VI increment-control condition.  The frequency
surrogates $\nu_n = \lambda_n / (2\pi)$ then faithfully encode $\lambda_n$
as an effective continuum parameterization within the declared interface
regime.  No continuum ontology is asserted; the notation is shorthand only.
\qed
\end{proof}

\subsection{O-SPEC.CD: Continuum Response-Density Interface}\label{sec:spec-CD}

\begin{theorem}[\status{B}]\label{thm:spec-CD-discharged}
\textup{(Continuum Response-Density Interface --- O-SPEC.CD Discharged.)}
\textup{[Bridge: conditional on O-SPEC.CI discharged
(Thm.~\ref{thm:spec-CI-discharged}); and the certified discrete response
classes admitting asymptotically coherent accumulation behavior in the
Book~VI sense on the declared spectroscopy continuum interface regime.]}
Response-density notation $\rho(\nu)\,\mathrm{d}\nu$ is licensed as
faithful shorthand for the asymptotic accumulation behavior of the certified
discrete spectral response classes on the declared continuum interface
regime.
\end{theorem}

\begin{proof}
By O-SPEC.CI, the continuum frequency parameterization is licensed.
The Book~VI Continuum Bridge Schema applies a second time: the discrete
response classes $\{\mathsf{Resp}(X_n; P)\}$ accumulate in the
continuum-interface regime according to the asymptotic coherence
hypothesis.  The response-density notation $\rho(\nu)\,\mathrm{d}\nu$
is the faithful encoding of this accumulation, with no additional continuum
structure assumed.  Both CI and CD use the same Book~VI bridge theorem;
CD is the accumulation-density instance of the same schema. \qed
\end{proof}

\subsection{O-SPEC.END: Endpoint Packaging}\label{sec:spec-END}

\begin{theorem}[\status{C}]\label{thm:spec-END-discharged}
\textup{(Spectroscopy Endpoint --- O-SPEC.END Discharged.)}
\textup{[Conditional on O-SPEC.PR through O-SPEC.CD all discharged
(Thms.~\ref{thm:spec-PR-discharged}--\ref{thm:spec-CD-discharged}); the
ICDC schema (Thm.~\ref{thm:ICDC}, Book~V, [C]); and the BR source theorem
(Thm.~\ref{thm:BR-source}, Book~V, [C]).]}
The endpoint datum
\[
  E_{\mathrm{Spec}} \;=\;
  \bigl(\mathcal{C}_{\mathrm{obs}},\;
        \mathcal{O}_{\mathrm{Spec}},\;
        G_{\mathrm{Spec}},\;
        \mathcal{W}\bigr)
\]
is a certified lawful spectroscopy descendant closure claim on the declared
gauge-response spectroscopy regime:
\begin{enumerate}[label=\textup{(\roman*)}]
  \item branch-visible resonance and transition structure are certified
    (O-SPEC.PR, O-SPEC.TR);
  \item spectral response classes are invariant under undeclared
    reinterpretation (O-SPEC.NL);
  \item transition-accessible spectral packets are generator-compatible
    (O-SPEC.GC);
  \item continuum frequency and response-density notation are licensed
    on the declared interface regime (O-SPEC.CI, O-SPEC.CD);
  \item the entire package is sourced through $\mathbf{U}$ via the BR
    coercive-transport route and is ICDC-coherent.
\end{enumerate}
\end{theorem}

\begin{proof}
Items~(i)--(iv) follow directly from the named discharged obligations.
Item~(v): by O-SPEC.OP, $L_{\mathrm{Spec}}$ is the restriction of the
Book~III coercive operator, which the BR source theorem
(Thm.~\ref{thm:BR-source}) certifies as a lawful specialization of the
upstream coercive-transport form on $\mathbf{U}$.  The ICDC schema
(Thm.~\ref{thm:ICDC}) applies with spectral no-loss control as the
one-step inheritance datum: the gap-subcritical ICDC reading
($\Theta_n = 1$) gives the spectroscopy-local statement that certified
spectral structure persists under admissible enlargement.  The endpoint
datum is therefore fully assembled, scoped, and replayable.
\qed
\end{proof}

\begin{remark}[\status{R}]\label{rem:spec-gauge-response-complete}
\textbf{Gauge-response spectroscopy: discharge status.}
Under the hypotheses of
Thms.~\ref{thm:spec-PR-discharged}--\ref{thm:spec-END-discharged}, the
gauge-response spectroscopy program reaches its endpoint.  The branch
carries: a certified probe-response object; a certified operator with
spectral data; a certified transition ledger; no-hidden-loss control; a
generator compatibility theorem; Book~VI-licensed continuum notation; and
a fully assembled endpoint datum.  The remaining open obligation is
O-SPEC.MAT (matter coupling), which is not needed for the
gauge-response regime.
\end{remark}

\subsection{O-SPEC.MAT: Matter Coupling Frontier (Remains Open)}
\label{sec:spec-MAT}

\begin{openobligation}[\status{O}]\label{ob:spec-MAT-open}
\textup{(O-SPEC.MAT: Matter Coupling --- Remains Open.)}
The extension of spectroscopy from the gauge-response regime to a
matter-rich spectroscopy branch requires: (1) the SM matter sector to be
certified beyond its current conditional status
(SM Branch, \texttt{ob:SM-matter}); (2) a lawful matter coupling theorem
harmonizing the certified matter observable class with the spectroscopy
observable family $\mathcal{O}_{\mathrm{Spec}}$; and (3) an extended
endpoint theorem incorporating both gauge-response and matter-response
spectral data.  O-SPEC.MAT is not discharged here.  The gauge-response
spectroscopy program (O-SPEC.PR through O-SPEC.END) is complete
independently of this obligation.
\end{openobligation}

\end{document}
